\documentclass[12pt]{article}
\usepackage[T1]{fontenc}
\usepackage[utf8]{inputenc}
\usepackage{lmodern}
\usepackage{textcomp}
\usepackage{enumitem}
\usepackage{geometry}
\usepackage{graphicx}
\usepackage{amsmath, amssymb}
\geometry{margin=1in}
\begin{document}
\begin{center}
Constitutional Law - Undergraduate Level Assessment 9 \
Total Marks: 40 \
Answer all questions.
\end{center}

\section*{Multiple Choice (10 marks)}
\begin{enumerate}[label=\arabic*.]
\item Can countries rely on their domestic law as an excuse to violate their obligations under international law?
\begin{enumerate}[label=(\alph*)]
    \item Domestic law always prevails over international law
    \item Only customary international law prevails over domestic law
    \item Obligations under international law prevail over domestic law
    \item Constitutional obligations always prevail over obligations under international law
\end{enumerate}
\item Which statement best describes the nature and function of Kelsen's Grundnorm?
\begin{enumerate}[label=(\alph*)]
    \item The ultimate source of a legal system's morality.
    \item The rule that distinguishes norms from habits of obedience.
    \item The constitution of a state.
    \item A presupposition that facilitates our understanding of the legal system.
\end{enumerate}
\item Is piracy under international (jure gentium) law subject to universal jurisdiction?
\begin{enumerate}[label=(\alph*)]
    \item Piracy jure gentium is subject to flag State jurisdiction
    \item Piracy jure gentium is subject to universal jurisdiction
    \item Piracy jure gentium is subject to port State jurisdiction
    \item Piracy jure gentium is subject to nationality-based jurisdiction
\end{enumerate}
\item Posner denies the autonomy of law on two grounds. Name one of them.
\begin{enumerate}[label=(\alph*)]
    \item He denies that law develops independently of social and economic forces.
    \item He claims that law is economically immoral.
    \item He rejects a positivist account of law.
    \item He opposes any sociological analysis of the law.
\end{enumerate}
\item Who was an exponent of "natural law with a variable content"?
\begin{enumerate}[label=(\alph*)]
    \item John Rawls
    \item Stammler
    \item Jerome Hall
    \item John Finns
\end{enumerate}
\end{enumerate}

\section*{True / False (10 marks)}
\textit{Answer True or False}
\begin{enumerate}[label=\arabic*.]
\setcounter{enumi}{5}
\item True or False: Gemeinschaft is considered a type of civic society, whereas Gesellschaft is viewed as a form of civil society.
\item True or False: In Hohfeld's scheme of jural relations, power and liability are indeed contradicting concepts.
\item True or False: Holland argued that law is solely a product of moral principles rather than state enforcement.
\item True or False: Dworkin suggests that equality of resources is measured when no one would prefer another's bundle of resources to their own.
\item True or False: The rule of recognition confers power and imposes duties on judges to decide cases within a legal system.
\end{enumerate}

\section*{Long Form Response (20 marks)}
\textit{Answer each question with a clear and complete explanation.}
\begin{enumerate}[label=\arabic*.]
\setcounter{enumi}{10}
\item Discuss the nature of the UK Constitution, focusing on its uncodified status and the various sources from which it derives. Analyze the implications of this structure for governance and legal interpretation.
\item Discuss Durkheim's concept of mechanical and organic solidarity and analyze their respective relationships to the types of law within society, focusing on repressive and restitutive law.
\end{enumerate}

\vfill
\noindent\textit{End of Paper}
\end{document}